
\begin{proposition}[A complete nodal-family computation]\label{cdp:st:B1}
Let $0<\varepsilon<4/27$, let $\Delta=\{|t|<\varepsilon\}$, and define
\[
 \mathcal C=\{y^2z=x^2(x+z)+tz^3\}\subset\PP^2\times\Delta,
 \qquad X=\mathcal C\times\PP^1.
\]
The projection $X\to\Delta$ is proper, its total space is smooth, and its central fibre is $S=C_0\times\PP^1$, where $C_0$ is an irreducible nodal cubic. For each $\lambda\in\C^*$ let $A_\lambda$ be the line bundle whose pullback to the normalization $\nu:\PP^1_u\to C_0$ is trivial and whose sections in this trivialization satisfy $h(1)=\lambda h(-1)$. Put $A=\operatorname{pr}_1^*A_\lambda$ and $d_\lambda=1$ for $\lambda=1$, and $d_\lambda=0$ otherwise. In the order $q=0,1,2$, the coherent cohomology dimensions are
\[
\begin{array}{c|c}
\text{sheaf on }S & (h^0,h^1,h^2)\\\hline
 N^*_{S/X}\otimes A &(d_\lambda,d_\lambda,0)\\
 K_S\otimes A &(1,0,0)\\
 T_S\otimes A &(1,0,0)\\
 \widetilde\Omega^1_S\otimes A &(0,1+d_\lambda,d_\lambda)\\
 \Omega^1_S\otimes A &(1,1+d_\lambda,d_\lambda)\\
 \Omega^1_X|_S\otimes A &(1,1+d_\lambda,d_\lambda).
\end{array}                                                    \tag{B.1}\label{cdp:eq:B.1}
\]
Here $T_S$ denotes the torsion of the cotangent sheaf, and $K_S$ its inverse image in the ambient cotangent sheaf. In particular the two proposed vanishings hold exactly for $\lambda\ne1$, and the ambient space of sections has dimension one for every $\lambda\in\C^*$. After shrinking $\Delta$, every $A_\lambda$ extends to a line bundle on $\mathcal C$ of the form $\OO_{\mathcal C}(P-Q)$ for holomorphic sections $P,Q$.
\end{proposition}

\begin{proof}
In the affine chart $z=1$, the defining equation is $F=y^2-x^3-x^2-t=0$. Its partial derivative in $t$ is $-1$, so the total space is smooth there. Its only point with $z=0$ is $[0:1:0]$. In the chart $y=1$ the equation is $z-x^3-x^2z-tz^3=0$, whose $z$ derivative is one at this point. Fibre singularities in the affine chart satisfy $y=0$ and $x(3x+2)=0$; they occur at $t=0,-4/27$. Our disc contains only the first. The quadratic term $y^2-x^2$ at the origin has two distinct linear factors, and its smoothing parameter has nonzero differential. The holomorphic Morse lemma with parameter therefore supplies local coordinates with $t=vw$. Properness follows from the closed inclusion into $\PP^2\times\Delta$.

The normalization is
\[
 u\longmapsto(x,y)=(u^2-1,u(u^2-1)),
\]
completed by the point at infinity. The two points above the node are $1,-1$. The descent sequence, an exact sequence of coherent $\OO_{C_0}$-modules, is
\[
0\longrightarrow A_\lambda\longrightarrow\nu_*\OO_{\PP^1}
 \xrightarrow{\ h\mapsto h(1)-\lambda h(-1)\ }\C_{\rm node}
 \longrightarrow0.                                      \tag{B.2}\label{cdp:eq:B.2}
\]
Its cohomology is the two-term complex $\C\xrightarrow{1-\lambda}\C$. Thus $h^0(A_\lambda)=h^1(A_\lambda)=d_\lambda$.

Write $V=\Omega^1_{\mathcal C}|_{C_0}$. The normalization map induces a surjective $\OO_{C_0}$-linear morphism $\Omega^1_{C_0}\to\nu_*\Omega^1_{\PP^1}$; at a node $vw=0$, it sends $dv$ to $(dv,0)$ and $dw$ to $(0,dw)$. Its kernel is the length-one module generated by $v\,dw=-w\,dv$. Consequently
\[
 T_{C_0}=\C_{\rm node},\qquad
 \widetilde\Omega^1_{C_0}=\nu_*\OO_{\PP^1}(-2).
\]
The kernel of $V\to\nu_*\Omega^1_{\PP^1}$ is $\nu_*\OO_{\PP^1}$. Locally its generators are $w\,dv$ and $v\,dw$ on the two branches; the identification is the normal-crossing conormal identification of Part A with the globally fixed conormal $dt$. Tensoring by $A_\lambda$ preserves exactness, and finite pushforward and the projection formula identify the two resulting normalization terms with $\nu_*\OO$ and $\nu_*\OO(-2)$. Hence
\[
 (h^0,h^1)(C_0,V\otimes A_\lambda)=(1,1),\qquad
 (h^0,h^1)(C_0,\Omega^1_{C_0}\otimes A_\lambda)=(1,1).
\]
The product morphisms give the direct decompositions
\[
 \Omega^1_X|_S\otimes A
 =(V\otimes A_\lambda)\boxtimes\OO
 \ \oplus\ A_\lambda\boxtimes\OO(-2),
\]
and the analogous decomposition for $\Omega^1_S$. The torsion-free quotient is obtained by replacing the first factor by $\nu_*\OO(-2)$. The torsion is $\C_{\rm node}\boxtimes\OO$, its ambient inverse image is $\nu_*\OO\boxtimes\OO$, and $N^*_{S/X}=\OO_S$. The coherent Kunneth isomorphism and $H^0(\PP^1,\OO)=\C$, $H^1(\PP^1,\OO(-2))=\C$ now give every entry of \eqref{cdp:eq:B.1}. All other cohomology of these two line bundles on $\PP^1$ vanishes.

We give the extension explicitly. Fix $Q(0)$ at $u=0$, and take $P(0)$ at $u=a=(1-\lambda)/(1+\lambda)$; for $\lambda=-1$ take $a=\infty$. For $\lambda\ne0$ both points are in the smooth locus. With $b=0$, the rational function $r(u)=(u-a)/(u-b)$ has divisor $P(0)-Q(0)$, and a local section in its pullback trivialization descends precisely with scalar
\[
 \frac{r(1)}{r(-1)}
 =\frac{(1-a)(1+b)}{(1-b)(1+a)}=\lambda.
\]
For $a=\infty$ the same expression is its limit, namely $-1$. The projection $\mathcal C\to\Delta$ is a submersion at these points, so the holomorphic implicit-function theorem extends them to sections on a smaller disc. The point at infinity itself is a constant section. For $\lambda=1$ take $P=Q$ identically. Sections through the smooth fibre locus are effective Cartier divisors in the smooth surface $\mathcal C$; their difference therefore defines the claimed extension. Pullback along $X\to\mathcal C$ extends $A$. This proves the claim also for all roots of unity, including the identity.
\end{proof}

The calculation recovers the local example in the supplied correction note, Lemma 4.3 and Theorem 4.4, with the gluing convention fixed by \eqref{cdp:eq:B.2}. The computation applies to roots of unity as well: the only exceptional scalar for the conormal vanishing is $1$. The symbolic computation and the exact gluing matrices are recorded in \texttt{code/B\_examples.py} and its output.


